\chapter{放送の社会的責任}

\section{放送の公共性と法的基盤}
\subsection{放送の公共性}

放送は公共財であり、単なる娯楽や情報提供ではない。そのため、\emp{社会に与える影響が大きい}特性を持つ。
電波は有限であり、すべての人が自由に利用できるわけではない。
この希少性のため、放送局は\emp{公共性}を担保する義務を負う。

公共性とは、放送が個人や企業の利益だけでなく、国民全体の知る権利や民主的社会の維持を目的とすることを意味する。
例えば、選挙時の報道では、有権者が公正な判断を下せる情報を提供する必要がある。

総務省によると、放送の目的は「国民の知る権利を保障し、健全な民主主義の発達に寄与すること」と規定されており、
これは\emp{放送法第1条}に明記されている（総務省, 2020）。これは単なる理念ではなく、放送局の運営や番組制作の判断基準として機能する。

\subsection{法的枠組み}
\begin{itemize}
  \item \emp{憲法21条（表現の自由）}：放送も含めた表現活動の自由を保障する。
  \item \emp{放送法第1条（放送目的）}：公共性・民主主義への寄与を明記。
  \item \emp{免許制と監督制度}：電波利用の公平性と秩序を維持するため、総務省による免許・監督制度がある。
\end{itemize}

憲法21条の表現の自由は、個人に広く認められる権利であるが、
放送は電波の有限性や公共性の観点から、一定の規律のもとに運用される必要がある。
したがって、放送現場における編集判断や番組構成は、
自由な表現と公共性の均衡を意識して行われなければならない。
よって、放送局は\emp{自由な表現}と\emp{公共的義務}の二重制約下で番組を制作する。
この制約は、政治的コメントの編集や災害報道の情報選択など、日々の現場判断に直結する。

\subsection{現場での判断}
各放送局が判断基準としているものの内容は、各放送局や連盟が定める\emp{放送基準}や\emp{番組基準}にも記載されている。
放送法第5条に基づき、各放送局は自ら番組基準を定め、それを公表することが義務付けられている。
これによって、国家権力による直接的な介入ではなく、放送界が自律的に表現の質を保つ仕組み(自律)をとっている。
例えば、日本民間放送連盟(民放連)の放送基準{\small(参考資料1参照)}では、「民主主義の精神にしたがい、基本的人権と世論を尊び、言論および表現の自由をまもり、法と秩序を尊重して社会の信頼にこたえる。」(日本民間放送連盟 放送基準から一部抜粋)とし、その後留意点が明記されていることが確認できる。

\section{公平・公正・政治的中立}
\subsection{}

\section{報道倫理とBPO}

\section{災害放送・社会的責任}

\section{名誉毀損・プライバシー・差別問題}